% ============================================================
\chapter{Discrete-Time Markov Chains}
\label{ch:discrete_markov}
% ============================================================

``The future depends on the past only through the present.'' This disarmingly simple statement captures the essence of the Markov property, formulated by Andrei Markov in 1906 while studying sequences of letters in Pushkin's poem \emph{Eugene Onegin}. Markov wanted to prove that the law of large numbers could apply beyond independent variables, and in doing so he discovered one of the most universal structures in applied mathematics. From search engines to genetic models, from queuing theory to MCMC algorithms, Markov chains are everywhere.

This chapter develops the theory of Markov chains with countable state spaces and discrete time: classification of states, recurrence, ergodicity, and convergence to equilibrium.

\section{The Markov Property}

\begin{definition}[Markov Chain]
Let $E$ be a countable set (the \emph{state space}). A process $(X_n)_{n \geq 0}$
with values in $E$, adapted to a filtration $(\mathcal{F}_n)_{n \geq 0}$, is a
\emph{Markov chain} if for every $n \geq 0$ and every $j \in E$:
\[
  \PP(X_{n+1} = j \mid \mathcal{F}_n) = \PP(X_{n+1} = j \mid X_n) \quad \text{a.s.}
\]
\end{definition}

\begin{intuition}[title=\textbf{Memorylessness}]
The Markov property means that the future of the process, given the present, is
independent of the past. All relevant past information is summarised in the
current state $X_n$.
\end{intuition}

\begin{definition}[Time Homogeneity]
A Markov chain is \emph{(time-)homogeneous} if the transition probabilities do
not depend on $n$: there exists a \emph{transition matrix} $P = (p_{ij})_{i,j \in E}$
with
\[
  p_{ij} = \PP(X_{n+1} = j \mid X_n = i), \quad \forall n \geq 0.
\]
The entries satisfy $p_{ij} \geq 0$ and $\sum_{j \in E} p_{ij} = 1$ for all $i$
(stochastic matrix).
\end{definition}

\begin{definition}[$n$-Step Transition Matrix]
The $n$-step transition matrix is $P^n = (p_{ij}^{(n)})$ where
\[
  p_{ij}^{(n)} = \PP(X_n = j \mid X_0 = i).
\]
\end{definition}

\begin{theorem}[Chapman--Kolmogorov Equations]
For all $m, n \geq 0$ and $i, j \in E$:
\[
  p_{ij}^{(m+n)} = \sum_{k \in E} p_{ik}^{(m)} p_{kj}^{(n)}.
\]
In matrix notation: $P^{m+n} = P^m P^n$.
\end{theorem}

\section{Strong Markov Property}

\begin{theorem}[Strong Markov Property]
\label{thm:strong_markov}
Let $(X_n)_{n \geq 0}$ be a homogeneous Markov chain and let $\tau$ be a stopping
time with $\PP(\tau < \infty) > 0$. Then, conditionally on $\{\tau < \infty\}$
and $X_{\tau} = i$, the process $(X_{\tau + n})_{n \geq 0}$ is a Markov chain
with the same transition matrix $P$, starting from $i$, independent of
$\mathcal{F}_{\tau}$.
\end{theorem}

\section{State Classification}

\begin{definition}[Communication]
We say $i$ \emph{leads to} $j$, written $i \to j$, if there exists $n \geq 1$
with $p_{ij}^{(n)} > 0$. We say $i$ and $j$ \emph{communicate}, written
$i \leftrightarrow j$, if $i \to j$ and $j \to i$.
\end{definition}

\begin{proposition}
The relation $\leftrightarrow$ is an equivalence relation on $E$. The equivalence
classes are called \emph{communicating classes}.
\end{proposition}

\begin{definition}[Irreducible Chain]
A Markov chain is \emph{irreducible} if $E$ forms a single communicating class,
i.e.\ $i \leftrightarrow j$ for all $i, j \in E$.
\end{definition}

\begin{definition}[Period]
The \emph{period} of a state $i$ is
\[
  d(i) = \gcd\{n \geq 1 : p_{ii}^{(n)} > 0\}.
\]
If $d(i) = 1$, the state $i$ is called \emph{aperiodic}. If the chain is
irreducible, all states have the same period $d$.
\end{definition}

\begin{definition}[Recurrence and Transience]
\label{def:recurrence}
Let $\tau_i = \inf\{n \geq 1 : X_n = i\}$ be the first return time to $i$.
\begin{itemize}
  \item State $i$ is \emph{recurrent} if $\PP_i(\tau_i < \infty) = 1$.
  \item State $i$ is \emph{transient} if $\PP_i(\tau_i < \infty) < 1$.
\end{itemize}
\end{definition}

\begin{theorem}[Recurrence Criterion]
\label{thm:recurrence_criterion}
A state $i$ is recurrent if and only if
\[
  \sum_{n=0}^{\infty} p_{ii}^{(n)} = +\infty.
\]
Equivalently, $i$ is transient if and only if $\sum_{n=0}^{\infty} p_{ii}^{(n)} < \infty$.
\end{theorem}

\begin{proof}
Let $N_i = \sum_{n=1}^{\infty} \mathbf{1}_{\{X_n = i\}}$ be the number of visits
to $i$ after time $0$. By the strong Markov property,
$\PP_i(N_i \geq k) = (\PP_i(\tau_i < \infty))^k$ for all $k \geq 1$. Hence:
\[
  \E_i[N_i] = \sum_{k=1}^{\infty} \PP_i(N_i \geq k)
  = \frac{\PP_i(\tau_i < \infty)}{1 - \PP_i(\tau_i < \infty)}
\]
(with the convention $= +\infty$ if $\PP_i(\tau_i < \infty) = 1$).
Since $\E_i[N_i] = \sum_{n=1}^{\infty} p_{ii}^{(n)}$, the result follows.
\end{proof}

\begin{definition}[Positive and Null Recurrence]
A recurrent state $i$ is:
\begin{itemize}
  \item \emph{positive recurrent} if $\E_i[\tau_i] < \infty$,
  \item \emph{null recurrent} if $\E_i[\tau_i] = +\infty$.
\end{itemize}
\end{definition}

\begin{keyformulas}[title=\textbf{State Classification}]
\[
\text{State } i \;\begin{cases}
  \text{transient} & \PP_i(\tau_i < \infty) < 1 \\
  \text{recurrent} & \PP_i(\tau_i < \infty) = 1
    \begin{cases}
      \text{null recurrent} & \E_i[\tau_i] = +\infty \\
      \text{positive recurrent} & \E_i[\tau_i] < +\infty
    \end{cases}
\end{cases}
\]
In an irreducible chain, all states have the same type.
\end{keyformulas}

\section{Stationary Measures}

\begin{definition}[Stationary Measure]
A measure $\pi = (\pi_j)_{j \in E}$ on $E$ (with $\pi_j \geq 0$) is
\emph{stationary} (or \emph{invariant}) for $P$ if:
\[
  \pi_j = \sum_{i \in E} \pi_i p_{ij}, \quad \forall j \in E,
\]
i.e.\ $\pi P = \pi$ in vector notation. If additionally $\sum_j \pi_j = 1$,
it is called a \emph{stationary distribution}.
\end{definition}

\begin{theorem}[Existence and Uniqueness]
\label{thm:stat_uniq}
Let $(X_n)$ be an irreducible Markov chain.
\begin{enumerate}
  \item A stationary measure exists if and only if the chain is recurrent.
  \item The stationary measure is unique (up to a multiplicative constant) and
    given by $\pi_j = 1/\E_j[\tau_j]$. It is normalisable (hence a distribution)
    if and only if the chain is positive recurrent.
\end{enumerate}
\end{theorem}

\section{Ergodic Theorem}

\begin{theorem}[Ergodic Theorem for Markov Chains]
\label{thm:ergodic}
Let $(X_n)_{n \geq 0}$ be an irreducible, aperiodic, positive recurrent Markov
chain with stationary distribution $\pi$. Then for all $i \in E$ and every
function $f : E \to \R$ with $\sum_j \abs{f(j)} \pi_j < \infty$:
\begin{enumerate}
  \item \textbf{Ergodic convergence}:
    $\displaystyle \frac{1}{n} \sum_{k=0}^{n-1} f(X_k) \xrightarrow{\text{a.s.}}
    \sum_{j \in E} f(j) \pi_j$.
  \item \textbf{Convergence in distribution}: $p_{ij}^{(n)} \to \pi_j$ as $n \to \infty$.
\end{enumerate}
\end{theorem}

\begin{warning}[title=\textbf{Role of Aperiodicity}]
Aperiodicity is needed for the convergence $p_{ij}^{(n)} \to \pi_j$. If the chain
is periodic with period $d$, one only has Ces\`aro convergence:
$\frac{1}{n}\sum_{k=0}^{n-1} p_{ij}^{(k)} \to \pi_j$.
\end{warning}

\section{Reversibility}

\begin{definition}[Reversibility]
A Markov chain with matrix $P$ is \emph{reversible} with respect to a distribution
$\pi$ if the \emph{detailed balance equations} hold:
\[
  \pi_i p_{ij} = \pi_j p_{ji}, \quad \forall i, j \in E.
\]
\end{definition}

\begin{proposition}
If $\pi$ satisfies detailed balance, then $\pi$ is a stationary distribution.
\end{proposition}

\begin{proof}
Summing over $i$: $\sum_i \pi_i p_{ij} = \sum_i \pi_j p_{ji} = \pi_j \sum_i p_{ji} = \pi_j$.
\end{proof}

\begin{example}[Random Walk on a Graph]
The random walk on a finite undirected graph $G = (V, E)$ with
$p_{ij} = 1/\deg(i)$ if $(i,j) \in E$ is reversible with respect to
$\pi_i = \deg(i) / (2\abs{E})$.
\end{example}

\section{Fundamental Examples}

\begin{example}[Random Walk on $\Z$]
Let $(X_n)$ be the simple random walk on $\Z$ with $p = \PP(Y_n = 1)$ and
$q = 1 - p = \PP(Y_n = -1)$.
\begin{itemize}
  \item If $p = q = 1/2$: the chain is null recurrent (no stationary distribution).
  \item If $p \neq 1/2$: the chain is transient.
\end{itemize}
\end{example}

\begin{example}[Random Walk on $\Z^d$]
The simple symmetric random walk on $\Z^d$ is:
\begin{itemize}
  \item recurrent for $d = 1$ and $d = 2$ (P\'olya's theorem),
  \item transient for $d \geq 3$.
\end{itemize}
\end{example}

\begin{example}[Ehrenfest Chain]
Two urns contain a total of $N$ balls. At each step, a ball is chosen uniformly
at random and moved to the other urn. If $X_n$ is the number of balls in the
first urn:
\[
  p_{i,i+1} = 1 - \frac{i}{N}, \quad p_{i,i-1} = \frac{i}{N}.
\]
The stationary distribution is $\pi_i = \binom{N}{i} 2^{-N}$ (binomial).
The chain is reversible.
\end{example}

\section{Hitting Times and the Gambler's Ruin}

\begin{proposition}[Hitting Times]
For an irreducible chain, the mean hitting time $h_i = \E_i[\tau_j]$ (for
fixed $j \neq i$) satisfies the system:
\[
  h_i = 1 + \sum_{k \neq j} p_{ik} h_k, \quad i \neq j, \qquad h_j = 0.
\]
\end{proposition}

\begin{example}[Gambler's Ruin]
A gambler starts with $i$ euros and plays against a casino. At each round, he
wins \$1 with probability $p$ and loses \$1 with probability $q = 1-p$. The game
ends when the fortune reaches $0$ or $N$. The ruin probability is:
\[
  \PP_i(\text{ruin}) = \begin{cases}
    \displaystyle \frac{(q/p)^i - (q/p)^N}{1 - (q/p)^N} & \text{if } p \neq q, \\[6pt]
    1 - i/N & \text{if } p = q = 1/2.
  \end{cases}
\]
\end{example}

\section{Exercises}

\begin{exercise}
Let $(X_n)$ be a Markov chain on $E = \{1, 2, 3\}$ with transition matrix
\[
  P = \begin{pmatrix} 0 & 1/2 & 1/2 \\ 1/3 & 1/3 & 1/3 \\ 1/2 & 0 & 1/2 \end{pmatrix}.
\]
Determine whether the chain is irreducible, compute its period and stationary
distribution.
\end{exercise}

\begin{exercise}
Show that in an irreducible chain, if one state is recurrent then all states
are recurrent. \emph{Hint: use $i \leftrightarrow j$.}
\end{exercise}

\begin{exercise}
Prove P\'olya's theorem: the simple symmetric random walk on $\Z^2$ is recurrent.
\emph{Hint: use the recurrence criterion (Theorem~\ref{thm:recurrence_criterion})
and an asymptotic expansion of $p_{00}^{(2n)}$.}
\end{exercise}

\begin{exercise}
Let $(X_n)$ be the Markov chain on $\N$ with $p_{i,i+1} = p$ and $p_{i,0} = q = 1-p$
for $i \geq 0$. Find the stationary distribution (if it exists) and conditions
on $p$ for positive recurrence.
\end{exercise}

\begin{exercise}
Prove that if $(X_n)$ is a Markov chain reversible with respect to $\pi$, then
the time-reversed chain $(X_{N-n})_{0 \leq n \leq N}$ has the same law as
$(X_n)_{0 \leq n \leq N}$ when $X_0 \sim \pi$.
\end{exercise}
